\cleardoublepage
\chapter{Appendix / 부록}

\section*{Important Formulas and Definitions / 중요한 공식과 정의}

\begin{itemize}
  \item \textbf{Place Value / 자리값:} Value of a digit = digit $\times$ place value.  
        \kr{숫자의 위치에 따라 값이 정해진다. 예: $4275 = 4000 + 200 + 70 + 5$}
  \item \textbf{Fraction to Decimal / 분수를 소수로:} $\frac{a}{b} = a \div b$.  
        \kr{분수는 나눗셈으로 변환할 수 있다.}
  \item \textbf{Prime Number / 소수:} A number greater than $1$ with exactly two factors: $1$ and itself.  
        \kr{1과 자기 자신만 약수로 갖는 수이다.}
  \item \textbf{Composite Number / 합성수:} A number with more than two factors.  
        \kr{약수가 두 개보다 많은 수이다.}
  \item \textbf{GCF / 최대공약수:} Product of the lowest powers of common primes.  
        \kr{공통 소인수의 가장 낮은 지수를 곱한 값이다.}
  \item \textbf{LCM / 최소공배수:} Product of the highest powers of all primes present.  
        \kr{모든 소인수의 가장 높은 지수를 곱한 값이다.}
\end{itemize}

\section*{Further Reading and Resources / 추가 학습 자료}

\begin{itemize}
  \item "Elementary Mathematics" by Liping Ma
  \item "How to Be a Math Genius" by Mike Goldsmith
  \item Khan Academy: \url{https://www.khanacademy.org/}
  \item Purplemath: \url{https://www.purplemath.com/}
\end{itemize}

\newpage

\cleardoublepage
\section*{Summary of Key Concepts / \kr{핵심 개념 요약}}

% ===== Chapter 1 =====
\subsection*{Chapter 1: Numbers \& Operations / \kr{1단원: 수와 연산}}
\begin{itemize}
  \item \textbf{Place Value / \kr{자리값}}: Each digit’s value depends on its position.  
  Example: $4275 = 4000 + 200 + 70 + 5$.
  \item \textbf{Decimals and Fractions / \kr{소수와 분수}}: Both represent parts of a whole.  
  Example: $0.75 = \tfrac{3}{4}$.
  \item \textbf{Comparing and Ordering / \kr{수의 비교와 정렬}}: Use $>,<,=$; convert to decimals for easier comparison.  
  Example: $\tfrac{2}{5} < \tfrac{3}{7}$.
  \item \textbf{Factors and Multiples / \kr{약수와 배수}}: Factors divide evenly; multiples are products.  
  Example: factors of $12$ are $1,2,3,4,6,12$.
  \item \textbf{Prime vs.\ Composite / \kr{소수와 합성수}}: Primes have exactly two factors (1 and itself); composites have more.  
  Example: $2,3,5$ are prime; $4,6,8$ are composite.
  \item \textbf{GCF and LCM / \kr{최대공약수와 최소공배수}}:  
  GCF = largest common factor; LCM = smallest common multiple.  
  Example: for $12,18$, GCF $=6$, LCM $=36$.
\end{itemize}

% ===== Chapter 2 =====
\subsection*{Chapter 2: Algebra / \kr{2단원: 대수}}
\begin{itemize}
  \item \textbf{Variables \& Expressions / \kr{변수와 식}}: Variables represent numbers; expressions combine numbers, variables, and operations.  
  \item \textbf{Simplifying Expressions / \kr{식 단순화}}: Combine like terms and use the distributive property.  
  \item \textbf{One-Step Equations / \kr{일차(한 단계) 방정식}}: Solved with one operation; apply the same to both sides.  
  \item \textbf{Multi-Step Equations / \kr{다단계 방정식}}: Use distribution → combine like terms → isolate the variable.  
  \item \textbf{Inequalities / \kr{부등식}}: Show ranges; flip the inequality sign when multiplying/dividing by a negative.  
\end{itemize}

% ===== Chapter 3 =====
\subsection*{Chapter 3: Ratios, Rates, and Percents / \kr{3단원: 비, 율, 백분율}}
\begin{itemize}
  \item \textbf{Ratios / \kr{비}}: Compare two numbers or quantities (part–part, part–whole).  
  \item \textbf{Rates / \kr{율}}: Ratios with different units (e.g., speed, cost per item).  
  \item \textbf{Unit Rate / \kr{단위율}}: Rate with denominator $1$ (e.g., per item, per hour).  
  \item \textbf{Proportions / \kr{비례식}}: Equality of two ratios; solved by cross multiplication.  
  \item \textbf{Percents / \kr{백분율}}: Mean “per hundred”; easily converted with fractions/decimals.  
\end{itemize}

\cleardoublepage
\section*{Final Words / \kr{맺음말}}

Congratulations on completing this book!  
You’ve worked through numbers, algebra, ratios, percents, and more.  
That’s not easy—especially while learning in two languages.  

\blockquote{\kr{이 책을 끝까지 마친 여러분을 축하합니다! 수, 대수학, 비와 백분율 등 다양한 개념을 공부했습니다. 두 언어로 동시에 배우는 일은 쉽지 않지만, 여러분은 해냈습니다.}}

Remember, math is like a language.  
The more you practice, the more fluent you become.  
Mistakes are part of learning, so don’t be discouraged.  

\blockquote{\kr{수학은 하나의 언어와 같습니다. 연습할수록 더 유창해집니다. 실수는 배움의 과정이니 절대 낙심하지 마세요.}}

\subsection*{Next Steps / \kr{다음 단계}}
\begin{itemize}
  \item Keep practicing problems every day, even small ones.
  \item Review the vocabulary lists in this book regularly.
  \item Explore extra resources like Khan Academy or Purplemath.
  \item Challenge yourself with harder questions step by step.
\end{itemize}

\blockquote{\kr{매일 조금씩 문제를 풀고, 책 속 어휘를 반복 복습하며, 온라인 자료도 활용하세요. 그리고 한 단계씩 더 어려운 문제에 도전해 보세요.}}

\subsection*{A Personal Note / \kr{개인적인 메시지}}

We wrote this book because we once struggled too.  
If you ever feel stuck, remember—you are not alone.  
Step by step, you’ll grow stronger in both math and English.  
One day, you’ll be the one helping others.  

\blockquote{\kr{저희도 같은 어려움을 겪었기에 이 책을 만들었습니다. 힘들 때가 있어도 혼자가 아님을 기억하세요. 차근차근 연습하다 보면 수학과 영어 모두에서 성장할 것입니다. 언젠가는 여러분이 다른 사람을 도와줄 차례가 올 것입니다.}}

\vspace{1cm}
\begin{center}
\textit{Thank you for learning with us. Keep going—you’ve got this!}  
\kr{\textit{끝까지 함께해 주셔서 감사합니다. 계속 나아가세요—여러분은 할 수 있습니다!}}
\end{center}
